\documentclass[12pt]{article}
\usepackage[margin=1in]{geometry}
\usepackage{amsmath,amssymb,amsthm}
\usepackage{url}
\usepackage{enumitem}

\newtheorem{theorem}{Theorem}[section]
\newtheorem{lemma}[theorem]{Lemma}
\newtheorem{proposition}[theorem]{Proposition}
\newtheorem{corollary}[theorem]{Corollary}
\theoremstyle{definition}
\newtheorem{definition}[theorem]{Definition}
\theoremstyle{remark}
\newtheorem{remark}[theorem]{Remark}

\numberwithin{equation}{section}

\newcommand{\Sph}{\mathbb S}
\newcommand{\Teich}{\operatorname{Teich}}
\newcommand{\Mod}{\operatorname{Mod}}
\newcommand{\Df}{\mathcal D_f}
\newcommand{\Lf}{L_f}
\newcommand{\Sf}{S_f}
\newcommand{\Ef}{E_f}
\newcommand{\Per}{\operatorname{Per}}
\newcommand{\HF}{\Phi_f}
\newcommand{\FAB}{\Phi_A^B}
\newcommand{\id}{\operatorname{id}}

\title{A Vanishing-Cycle Dehn Twist for the Hironaka--Koch
$\Per_4(0)$ Deformation Space: Preservation, Base-Component
Triviality, and the Open Infinite-Orbit Problem}
\author{}
\date{}

\begin{document}
\maketitle

\begin{abstract}
We give a partial answer to the problem of exhibiting an explicit simple
closed curve whose Dehn-twist power preserves the Hironaka--Koch
deformation space attached to a quadratic rational map with a
superattracting $4$-cycle and one free critical value, and whose action
on the connected components has an infinite orbit.  For the concrete
Hironaka--Koch map
\[
f(z)=\frac{(4z-3)(z+2)}{4z^{2}},
\]
with superattracting cycle
$0\mapsto\infty\mapsto1\mapsto\tfrac34\mapsto0$ and free critical value
$v=121/96$, let $\gamma_{0,v}$ be any essential simple closed curve
separating $\{0,v\}$ from $\{\infty,1,\tfrac34\}$.  We prove that the
square $T_{\gamma_{0,v}}^{2}$ of the Dehn twist about this curve
preserves the deformation space $\Df$.  We also prove, using the
base-point-component form of the Hironaka--Koch monodromy subgroup
$\Ef$, that $T_{\gamma_{0,v}}^{2}\in\Ef$, so it fixes the base
connected component of $\Df$.  Thus the explicit curve
$\gamma_{0,v}$ and the power $2$ solve the preservation half of the
problem, but they do not supply the requested infinite orbit.  The known
Hironaka--Koch infinite-component mechanism is produced by a braid
element that is not identified here with a Dehn twist about a single
simple closed curve.
\end{abstract}

\section{Introduction}
\label{sec:intro}

Let $f:\Sph^{2}\to\Sph^{2}$ be a branched cover, let $V(f)$ denote its
set of critical values, and let $A,B\subset\Sph^{2}$ be finite sets.
Whenever $f(A)\cup V(f)\subset B$, pulling back complex structures on
$(\Sph^{2},B)$ yields a holomorphic map
\[
f^{*}:\Teich(B)\longrightarrow\Teich(A),
\]
and whenever $A\subset B$, forgetting the points $B\setminus A$ yields
\[
i_{*}:\Teich(B)\longrightarrow\Teich(A).
\]
When $f(A)\cup V(f)\cup A\subset B$, the deformation space is the
equalizer
\[
\Df=\{\tau\in\Teich(B): f^{*}(\tau)=i_{*}(\tau)\}.
\]

The original problem asks for an explicit simple closed curve
$\gamma\subset(\Sph^{2},B)$ such that some power of the Dehn twist
$T_{\gamma}$ preserves $\Df$ and acts on the set of connected components
$\pi_{0}(\Df)$ with an infinite orbit.  We work with the specific
quadratic rational map used by Hironaka and Koch in \cite[\S4]{HK}:
\begin{equation}
\label{eq:HKmap}
f(z)=\frac{(4z-3)(z+2)}{4z^{2}}.
\end{equation}
Its critical points are $0$ and $12/5$.  The point $0$ lies in the
superattracting $4$-cycle
\[
A=\left\{0,\infty,1,\tfrac34\right\},\qquad
0\longmapsto\infty\longmapsto1\longmapsto\tfrac34\longmapsto0,
\]
while the second critical value is the free critical value
\[
v=f\!\left(\tfrac{12}{5}\right)=\frac{121}{96}.
\]
We set
\[
B=A\cup\{v\}.
\]

Let $\gamma_{0,v}$ be any essential simple closed curve in
$\Sph^{2}\setminus B$ bounding a disk whose intersection with $B$ is
exactly $\{0,v\}$: it separates $\{0,v\}$ from
$\{\infty,1,\tfrac34\}$.  Geometrically, $\gamma_{0,v}$ is the
vanishing cycle of the collision of the free critical value $v$ with the
periodic critical point $0$.

\medskip
The present paper proves the following two results, and records honestly
what is not proved.

\smallskip
\noindent\textbf{Theorem A (preservation).}
For each $a\in\{0,1,\tfrac34\}$, if $\gamma_{a,v}$ is any essential
simple closed curve separating $\{a,v\}$ from the other three points of
$B$, then $T_{\gamma_{a,v}}^{2}\in\Sf$, and consequently
\[
T_{\gamma_{a,v}}^{2}\cdot\Df=\Df.
\]
In particular, $\gamma=\gamma_{0,v}$ satisfies the preservation
requirement of the problem with power $2$.

\smallskip
\noindent\textbf{Theorem B (base-component triviality).}
With the Hironaka--Koch base point and marking conventions, the
vanishing-cycle twist satisfies
\[
T_{\gamma_{0,v}}^{2}\in\Ef,
\]
where $\Ef$ is the image of the monodromy of the Hironaka--Koch
parameter space $V=\Per_{4}(0)^{*}$.  Consequently, if $D_{0}$ is the
connected component of $\Df$ containing the base point $\tau_{*}$, then
$T_{\gamma_{0,v}}^{2}(D_{0})=D_{0}$.

Thus the same explicit curve $\gamma=\gamma_{0,v}$ that supplies the
preservation half of the problem does not supply an infinite orbit: its
square lies in $\Ef$ and fixes the base component.  The infinite-orbit
requirement for a Dehn twist about a single simple closed curve remains
open.  The known Hironaka--Koch infinite-component phenomenon is
produced by the braid element $g=s_{2}cr_{1}s_{1}cr_{2}$ of
\cite[Theorem~1.3 and Remark~4.7]{HK}, which we cite but do not identify
with a Dehn twist.
The paper is organised as follows.  Section~\ref{sec:prelim} recalls the
Hironaka--Koch framework, states their stabiliser criterion, and derives
the base-component lemma we use.  Section~\ref{sec:normal-form} proves
the normal-form algebra for quadratic rational maps with a
superattracting $4$-cycle and computes the multiplicity-two collision
between $v$ and the periodic critical point.
Section~\ref{sec:preservation} verifies the concrete data
of~\eqref{eq:HKmap}, establishes the structure of the preimage of the
vanishing cycle, and proves Theorem~A.  Section~\ref{sec:Ef} proves
Theorem~B.  Section~\ref{sec:open} summarises the consequences and
records the precise open points.

\section{Preliminaries and Hironaka--Koch input}
\label{sec:prelim}

Throughout, $\Sph^{2}$ is the Riemann sphere.  For a finite set
$P\subset\Sph^{2}$ with $|P|\ge3$, let $\Teich(P)$ denote the
Teichm\"uller space of the marked sphere $(\Sph^{2},P)$, and let
\[
\Mod_{P}=\pi_{0}\operatorname{Homeo}^{+}(\Sph^{2},P)
\]
denote the pure mapping class group: orientation-preserving
homeomorphisms fixing $P$ pointwise, modulo isotopy through such
homeomorphisms.  For an essential simple closed curve
$\gamma\subset\Sph^{2}\setminus P$, we write $T_{\gamma}$ for the right
Dehn twist about $\gamma$.

If $Q\subset P$ has cardinality $2$, we let $\gamma_{Q}$ denote any
essential simple closed curve in $\Sph^{2}\setminus P$ bounding a disk
$D$ with $D\cap P=Q$.  Such a curve separates $\Sph^{2}$ into a disk
containing $Q$ and a disk containing $P\setminus Q$.  For the concrete
map of Section~\ref{sec:preservation}, we use the notation
\[
A=\{0,\infty,1,\tfrac34\},\qquad
B=A\cup\{v\},\qquad
v=\tfrac{121}{96},
\]
and for $a\in A$ we write $\gamma_{a,v}$ for $\gamma_{\{a,v\}}$.

\subsection{The Hironaka--Koch deformation setup}

We work inside the framework of Hironaka--Koch \cite[§2.5--§2.6]{HK}.
A mapping class $g\in\Mod_{B}$ is called \emph{$f$-liftable} if there
exists a class $\HF(g)\in\Mod_{A}$ such that
\[
f^{*}(g\cdot\tau)=\HF(g)\cdot f^{*}(\tau)\qquad\text{for all }\tau\in\Teich(B).
\]
Equivalently, $g$ admits a lift to a homeomorphism of
$(\Sph^{2},f^{-1}(B))$ fixing $A$ pointwise after forgetting
$f^{-1}(B)\setminus A$.  The liftable classes form a subgroup
$\Lf\subset\Mod_{B}$, and the lifting homomorphism is
\[
\HF:\Lf\longrightarrow\Mod_{A}.
\]
Because $A\subset B$, forgetting $B\setminus A$ gives a homomorphism
\[
\FAB:\Mod_{B}\longrightarrow\Mod_{A}.
\]
The \emph{special liftable} subgroup is the equalizer
\[
\Sf=\{g\in\Lf:\HF(g)=\FAB(g)\}.
\]

We use the following standard Hironaka--Koch criterion.  We cite it
rather than reprove it, and we check its hypotheses for our concrete
data in Section~\ref{sec:preservation}.

% RECORDS: 564,262
\begin{proposition}[Hironaka--Koch {\cite[Proposition~2.1]{HK}}]
\label{prop:HK21}
Let $f:\Sph^{2}\to\Sph^{2}$ be a branched cover and let
$A,B\subset\Sph^{2}$ be finite sets with $A\subset B$, $|A|\ge3$,
$|B|\ge3$, and
\[
f(A)\cup V(f)\cup A\subset B.
\]
With $\Lf,\HF,\FAB,\Sf$ as above, if $g\in\Lf$ then the following are
equivalent:
\begin{enumerate}[label=(\roman*)]
\item\label{it:stab-nonempty} $g\cdot\Df\cap\Df\neq\varnothing$;
\item\label{it:stab-eq} $g\cdot\Df=\Df$;
\item\label{it:stab-Sf} $g\in\Sf$.
\end{enumerate}
\end{proposition}
\begin{proof}
This is \cite[Proposition~2.1]{HK}.  The hypotheses used there are
exactly that both Teichm\"uller spaces are non-exceptional
($|A|,|B|\ge3$) so that $f^{*}$ and $i_{*}$ are well defined, and that
$f(A)\cup V(f)\cup A\subset B$ so that the deformation space
$\Df=\{\tau\in\Teich(B):f^{*}(\tau)=i_{*}(\tau)\}$ is meaningful.

For the map $f$ of~\eqref{eq:HKmap}, Lemma~\ref{lem:HK-map} below
computes $f(A)=A$ and $V(f)=\{\infty,v\}$, so
\[
f(A)\cup V(f)\cup A=A\cup\{v\}=B,
\]
and $|A|=4$, $|B|=5$.  Hence the hypotheses of the cited proposition
are satisfied.
\end{proof}

\subsection{The Hironaka--Koch parameter space and the subgroup $\Ef$}

Let
\[
V=\Per_{4}(0)^{*}
\]
be the Hironaka--Koch parameter space of quadratic rational maps with a
marked superattracting $4$-cycle and one free critical value, in the
normalization of \cite[\S4]{HK}.  Let $*_{V}\in V$ be the base point
corresponding to the concrete map~\eqref{eq:HKmap}, and let
$\tau_{*}\in\Teich(B)$ be the associated marked point.  Hironaka--Koch
construct a monodromy homomorphism
\[
\rho:\pi_{1}(V,*_{V})\longrightarrow\Sf
\]
and define
\[
\Ef=\rho(\pi_{1}(V,*_{V})).
\]
In \cite{HK} this subgroup is written $i_{*}(\pi_{1}(V,*_{V}))$; to
avoid a clash of notation with the forgetful map
$i_{*}:\Teich(B)\to\Teich(A)$, we write $\rho$ for the monodromy
homomorphism.

We shall not use the full coset description of the connected components
of $\Df$ from \cite[Proposition~3.4]{HK}.  All we need is the following
base-point-component consequence.

% RECORDS: 643
\begin{lemma}
\label{lem:base-component}
Let $D_{0}$ be the connected component of $\Df$ containing the base
point $\tau_{*}$.  If $g\in\Ef$, then $g(D_{0})=D_{0}$.
\end{lemma}
\begin{proof}
Choose a loop $\ell$ in $V$ based at $*_{V}$ with $\rho([\ell])=g$.
By the Hironaka--Koch construction of the monodromy action, the
basepoint marking transported along the loop gives a path of marked
structures on which the deformation equation holds at every parameter;
thereore the lifted path lies in $\Df$.  Starting at $\tau_{*}$ the
endpoint of this path is $g\cdot\tau_{*}$ (in the action convention
of~\cite{HK}), so $g\cdot\tau_{*}\in D_{0}$.

Since $\Ef\subset\Sf$, Proposition~\ref{prop:HK21} gives
$g\cdot\Df=\Df$.  Hence $g(D_{0})$ is a connected component of $\Df$.
This component contains $g\cdot\tau_{*}$, and $D_{0}$ also contains
$g\cdot\tau_{*}$.  Since distinct components are disjoint,
$g(D_{0})=D_{0}$.
\end{proof}

\begin{remark}
The conclusion of Lemma~\ref{lem:base-component} is deliberately
limited to the base component.  If $\Ef$ were known to be normal in
$\Sf$, one could draw further conclusions about the action on all
components; the arguments of this paper do not establish normality, and
we do not rely on it.
\end{remark}

\section{Normal form and the double collision}
\label{sec:normal-form}

The next theorem records the one-parameter normal form for quadratic
rational maps with a superattracting $4$-cycle, computes the second
critical point and its value, and identifies the collision between the
free critical value and the periodic critical point together with its
local multiplicity.

% RECORDS: 262
\begin{theorem}
\label{thm:normal-form}
For $\alpha$ away from the exceptional loci on which the formulas below
degenerate, define
\[
f_{\alpha}(z)=\frac{\alpha(z-\alpha)^{2}}{(z-1)(z-\alpha^{3})}.
\]
Then:
\begin{enumerate}[label=(\roman*)]
\item\label{it:nf-cycle} $\alpha$ is a critical point, and
\[
\alpha\longmapsto0\longmapsto1\longmapsto\infty\longmapsto\alpha
\]
is a superattracting $4$-cycle.
\item\label{it:nf-c2} The second critical point is
\[
c_{2}(\alpha)=-\frac{\alpha(\alpha^{2}-\alpha-1)}{\alpha^{2}+\alpha-1}.
\]
\item\label{it:nf-v} Its critical value is
\[
v(\alpha)=f_{\alpha}(c_{2}(\alpha))
=\frac{4\alpha^{2}(\alpha+1)}{(\alpha^{2}+\alpha+1)^{2}}.
\]
\item\label{it:nf-collision} The difference of the free critical value
from the periodic critical point factors as
\[
v(\alpha)-\alpha
=-\frac{\alpha(\alpha^{2}+\alpha-1)^{2}}{(\alpha^{2}+\alpha+1)^{2}}.
\]
Hence $v(\alpha)=\alpha$ exactly when $\alpha^{2}+\alpha-1=0$, and at
such parameters the collision has multiplicity $2$; additionally
$c_{2}(\alpha)\to\infty=f_{\alpha}^{3}(\alpha)$.
\item\label{it:nf-mono} If $\gamma_{\alpha}$ is any simple closed
curve separating $\{\alpha,v(\alpha)\}$ from $\{0,1,\infty\}$, then the
local monodromy of the family $f_{\alpha}$ around the collision
parameter is $T_{\gamma_{\alpha}}^{2}$, up to the usual orientation
convention for Dehn twists.
\end{enumerate}
\end{theorem}
\begin{proof}
Set
\[
N(z)=\alpha(z-\alpha)^{2},\qquad
D(z)=(z-1)(z-\alpha^{3}),
\]
so that $f_{\alpha}=N/D$.

\emph{(i).}  The numerator has a double zero at $z=\alpha$, so $f$
attains the value $0$ there to local degree $2$, and $\alpha$ is a
critical point.  Direct evaluation gives
\[
f_{\alpha}(0)=\frac{\alpha(-\alpha)^{2}}{(-1)(-\alpha^{3})}=1,\qquad
f_{\alpha}(\infty)=\alpha,
\]
while $D(1)=0$ and $N(1)\neq0$ away from exceptional $\alpha$, so
$f_{\alpha}(1)=\infty$.  Hence
\[
\alpha\longmapsto0\longmapsto1\longmapsto\infty\longmapsto\alpha
\]
is a $4$-cycle, and it is superattracting because $\alpha$ is critical.

\emph{(ii).}  Differentiating,
\[
f_{\alpha}'(z)=\frac{N'(z)D(z)-N(z)D'(z)}{D(z)^{2}},
\qquad N'(z)=2\alpha(z-\alpha),
\]
so
\[
f_{\alpha}'(z)
=\frac{\alpha(z-\alpha)\bigl(2D(z)-(z-\alpha)D'(z)\bigr)}{D(z)^{2}}.
\]
One critical point is $z=\alpha$.  Writing
$D(z)=z^{2}-(1+\alpha^{3})z+\alpha^{3}$ and
$D'(z)=2z-(1+\alpha^{3})$, we compute
\[
\begin{aligned}
2D(z)-(z-\alpha)D'(z)
&=2\bigl(z^{2}-(1+\alpha^{3})z+\alpha^{3}\bigr)
 -(z-\alpha)\bigl(2z-(1+\alpha^{3})\bigr)\\
&=-(\alpha-1)\Bigl((\alpha^{2}+\alpha-1)z
   +\alpha(\alpha^{2}-\alpha-1)\Bigr).
\end{aligned}
\]
The remaining zero of this factor yields the second critical point
\[
c_{2}(\alpha)
=-\frac{\alpha(\alpha^{2}-\alpha-1)}{\alpha^{2}+\alpha-1}.
\]

\emph{(iii).}  Set $\beta=c_{2}(\alpha)$.  At the critical point, the
equation $2D(\beta)=(\beta-\alpha)D'(\beta)$ holds, hence
\[
v(\alpha)=f_{\alpha}(\beta)
=\frac{\alpha(\beta-\alpha)^{2}}{D(\beta)}
=\frac{2\alpha(\beta-\alpha)}{D'(\beta)}.
\]
Now
\[
\beta-\alpha
=-\frac{\alpha(\alpha^{2}-\alpha-1)}{\alpha^{2}+\alpha-1}-\alpha
=-\frac{2\alpha(\alpha^{2}-1)}{\alpha^{2}+\alpha-1},
\]
and
\[
D'(\beta)
=2\beta-(1+\alpha^{3})
=-\frac{(\alpha-1)(\alpha^{2}+\alpha+1)^{2}}{\alpha^{2}+\alpha-1}.
\]
Substituting,
\[
v(\alpha)
=\frac{2\alpha\cdot\left(-\dfrac{2\alpha(\alpha^{2}-1)}{\alpha^{2}+\alpha-1}\right)}
        {-\dfrac{(\alpha-1)(\alpha^{2}+\alpha+1)^{2}}{\alpha^{2}+\alpha-1}}
=\frac{4\alpha^{2}(\alpha+1)}{(\alpha^{2}+\alpha+1)^{2}}.
\]

\emph{(iv).}  Combining (iii) with the common denominator,
\[
\begin{aligned}
v(\alpha)-\alpha
&=\frac{4\alpha^{2}(\alpha+1)-\alpha(\alpha^{2}+\alpha+1)^{2}}
        {(\alpha^{2}+\alpha+1)^{2}}\\
&=\frac{-\alpha\bigl(\alpha^{4}+2\alpha^{3}-\alpha^{2}-2\alpha+1\bigr)}
        {(\alpha^{2}+\alpha+1)^{2}}\\
&=-\frac{\alpha(\alpha^{2}+\alpha-1)^{2}}
        {(\alpha^{2}+\alpha+1)^{2}}.
\end{aligned}
\]
Thus $v(\alpha)=\alpha$ if and only if $\alpha^{2}+\alpha-1=0$.  At
such $\alpha$, $\alpha^{2}+\alpha+1=2$, so the prefactor is nonzero,
and since $\alpha^{2}+\alpha-1$ has simple zeros, the collision has
multiplicity $2$.

At a root of $\alpha^{2}+\alpha-1$, the denominator
$\alpha^{2}+\alpha-1$ of $c_{2}(\alpha)$ vanishes.  Using
$\alpha^{2}=1-\alpha$,
\[
\alpha^{2}-\alpha-1=(1-\alpha)-\alpha-1=-2\alpha,
\]
so the numerator of $c_{2}$ tends to $2\alpha^{2}\neq0$ because
$\alpha\neq0$ on the collision locus; hence
$c_{2}(\alpha)\to\infty$.  Since $f_{\alpha}^{3}(\alpha)=\infty$ by
(i), the limiting second critical point coincides with the cycle point
$f_{\alpha}^{3}(\alpha)$.

\emph{(v).}  Let $\alpha_{0}$ be a root of $\alpha^{2}+\alpha-1$ and
write $\alpha-\alpha_{0}=t$.  By (iv),
\[
v(\alpha)-\alpha =u(t)\,t^{2},\qquad u(0)\neq0.
\]
Along the loop $t=\varepsilon e^{i\theta}$, $0\le\theta\le2\pi$, the
relative coordinate $v(\alpha)-\alpha$ winds twice around $0$.  Choose a
closed disk $U$ containing $\alpha$ and $v(\alpha)$ and disjoint from
$\{0,1,\infty\}$.  Inside $U$, the two marked points move as the pure
braid with relative winding number two.  Under the standard
braid-to-mapping-class dictionary, in the twice-punctured disk a pure
braid with relative winding number one represents the Dehn twist about
the simple closed curve that surrounds both punctures, and relative
winding number two represents its square.  See Farb--Margalit
\cite[Chapter~9, Theorem~9.1 and \S9.1]{FM} for the identification of
the pure braid group on a punctured disk with a subgroup of the mapping
class group.  Extending by the identity outside $U$ gives the local
monodromy $T_{\gamma_{\alpha}}^{2}$, up to orientation.
\end{proof}

\begin{remark}
\label{rem:coord-change}
The normalization of Hironaka--Koch \cite[\S4]{HK} places the periodic
critical point at the value $0$ by a M\"obius change of coordinates.  In
that coordinate, the curve $\gamma_{\alpha}$ of
Theorem~\ref{thm:normal-form} becomes the curve $\gamma_{0,v}$ that
separates $\{0,v\}$ from the other three points of $B$.
\end{remark}

\section{The concrete map and liftable squared twists}
\label{sec:preservation}

We now specialize to the map~\eqref{eq:HKmap},
\[
f(z)=\frac{(4z-3)(z+2)}{4z^{2}}
=\frac{4z^{2}+5z-6}{4z^{2}}.
\]
Its superattracting $4$-cycle is
\[
A=\left\{0,\infty,1,\tfrac34\right\},
\]
and its free critical value is
\[
v=\tfrac{121}{96},\qquad
B=A\cup\{v\}.
\]

% RECORDS: 564
\begin{lemma}
\label{lem:HK-map}
For the map $f$ of~\eqref{eq:HKmap}, the following hold.
\begin{enumerate}[label=(\roman*)]
\item\label{it:hk-cycle} The cycle is
\[
0\longmapsto\infty\longmapsto1\longmapsto\tfrac34\longmapsto0.
\]
\item\label{it:hk-crit} The critical points are $0$ and $12/5$, and the
branch values are $\infty$ and $v=121/96$.
\item\label{it:hk-deck} The nontrivial deck transformation of the
degree-two map is
\[
\tau(z)=\frac{6z}{5z-6},
\]
and $\tau^{2}=\id$ and $f\circ\tau=f$.
\item\label{it:hk-fibers} The fibers over $B$ are
\[
\begin{aligned}
f^{-1}(0)&=\left\{\tfrac34,\,-2\right\},\\
f^{-1}(\infty)&=\{0\}\quad\text{with branching},\\
f^{-1}(1)&=\left\{\tfrac65,\,\infty\right\},\\
f^{-1}\!\left(\tfrac34\right)&=\{1,\,-6\},\\
f^{-1}(v)&=\left\{\tfrac{12}{5}\right\}\quad\text{with branching}.
\end{aligned}
\]
\item\label{it:hk-factor} The map factors as
\[
f=\phi\circ q\circ m,\qquad
m(z)=\frac{5z}{5z-12},\quad q(w)=w^{2},\quad
\phi(s)=\frac{121s-25}{96s}.
\]
\end{enumerate}
\end{lemma}
\begin{proof}
\emph{(i).}  Evaluation gives
\[
f(0)=\infty,\qquad f(\infty)=1,\qquad
f(1)=\frac{4+5-6}{4}=\frac34,
\]
and
\[
f\!\left(\tfrac34\right)
=\frac{\tfrac94+\tfrac{15}{4}-6}{\tfrac94}=0,
\]
so the claimed $4$-cycle holds.

\emph{(ii).}  In the local coordinate $w=1/f$ at $f(0)=\infty$,
\[
w(z)=\frac{4z^{2}}{4z^{2}+5z-6},
\]
which has a zero of order $2$ at $z=0$.  Therefore $0$ has local degree
$2$: it is a critical point with critical value $\infty$.  Away from
$0$ we differentiate:
\[
f'(z)=\frac{(8z+5)4z^{2}-(4z^{2}+5z-6)8z}{16z^{4}}
=\frac{12-5z}{4z^{3}}.
\]
Hence the only finite critical point other than $0$ is $z=12/5$.  Its
value is
\[
f\!\left(\tfrac{12}{5}\right)
=\frac{\left(\tfrac{48-15}{5}\right)\left(\tfrac{12+10}{5}\right)}
       {4\cdot\tfrac{144}{25}}
=\frac{\tfrac{33}{5}\cdot\tfrac{22}{5}}{\tfrac{576}{25}}
=\frac{726}{576}=\frac{121}{96}=v.
\]
Finally, $\infty$ is not critical: with $u=1/z$ near $\infty$ and target
coordinate $W=f(1/u)-1$ near $f(\infty)=1$,
\[
W(u)=\frac{4/u^{2}+5/u-6}{4/u^{2}}-1
=\frac{5}{4}u-\frac{3}{2}u^{2},
\]
whose derivative at $u=0$ is $5/4\neq0$.  Therefore
$V(f)=\{\infty,v\}$.

\emph{(iii) and (v).}  First verify $\tau^{2}=\id$:
\[
\tau(\tau(z))=\frac{6\cdot\tfrac{6z}{5z-6}}{\tfrac{5\cdot6z}{5z-6}-6}
=\frac{\tfrac{36z}{5z-6}}{\tfrac{30z-6(5z-6)}{5z-6}}
=\frac{36z}{36}=z.
\]
Next, since
\[
4\tau^{2}+5\tau-6
=\frac{144z^{2}+180z-216}{(5z-6)^{2}},
\qquad
4\tau^{2}=\frac{144z^{2}}{(5z-6)^{2}},
\]
we have
\[
f(\tau(z))=\frac{144z^{2}+180z-216}{144z^{2}}
=1+\frac{5}{4z}-\frac{3}{2z^{2}}=f(z),
\]
which proves $f\circ\tau=f$.  For the factorization, if
$w=m(z)=5z/(5z-12)$, then $q(w)=w^{2}=25z^{2}/(5z-12)^{2}$ and
\[
\phi(q(m(z)))
=\frac{121\cdot\tfrac{25z^{2}}{(5z-12)^{2}}-25}
       {96\cdot\tfrac{25z^{2}}{(5z-12)^{2}}}
=\frac{25\bigl(121z^{2}-(5z-12)^{2}\bigr)}{2400z^{2}}.
\]
Since $(5z-12)^{2}=25z^{2}-120z+144$, the numerator is
$25(96z^{2}+120z-144)$, so
\[
\phi(q(m(z)))=\frac{2400z^{2}+3000z-3600}{2400z^{2}}
=1+\frac{5}{4z}-\frac{3}{2z^{2}}=f(z).
\]

\emph{(iv).}  We compute the fibers directly.
\[
f(z)=0\iff (4z-3)(z+2)=0\iff z\in\left\{\tfrac34,-2\right\},
\]
\[
f(z)=1\iff 4z^{2}+5z-6=4z^{2}\iff z=\tfrac65,
\]
and additionally $z=\infty$ satisfies $f(\infty)=1$, so
$f^{-1}(1)=\{\tfrac65,\infty\}$.  Next,
\[
f(z)=\tfrac34\iff 4(4z-3)(z+2)=12z^{2}
\iff z^{2}+5z-6=0\iff z\in\{1,-6\},
\]
so $f^{-1}(\tfrac34)=\{1,-6\}$.  For the value $v$, we have
\[
f(z)=v=\tfrac{121}{96}
\iff 384z^{2}+480z-576=484z^{2},
\]
hence
\[
100z^{2}-480z+576=0
\iff (5z-12)^{2}=0,
\]
so $f^{-1}(v)=\{12/5\}$ with multiplicity $2$.  Finally,
$f^{-1}(\infty)=\{0\}$ because $4z^{2}=0$ is the only solution of the
vanishing denominator, with multiplicity $2$.
\end{proof}

\begin{remark}
\label{rem:fminus-one-B}
From Lemma~\ref{lem:HK-map},
\[
f^{-1}(B)=\left\{\tfrac34,-2,\,0,\,\tfrac65,\,\infty,\,1,\,-6,\,\tfrac{12}{5}\right\},
\]
so that
\[
A\subset f^{-1}(B),\qquad
f^{-1}(B)\setminus A=\left\{-2,\,\tfrac65,\,-6,\,\tfrac{12}{5}\right\}.
\]
In particular, $f$ restricts to an unramified double cover
$f:\Sph^{2}\setminus f^{-1}(B)\to\Sph^{2}\setminus B$ with deck group
$\langle\tau\rangle$.
\end{remark}

\subsection{The preimage of a $v$-pair disk}

% RECORDS: 564
\begin{lemma}
\label{lem:disk-preimage}
Let $a\in\{0,1,\tfrac34\}$ and let $\gamma=\gamma_{a,v}$ be a simple
closed curve bounding a disk $D$ with $D\cap B=\{a,v\}$.  Then
$\delta=f^{-1}(\gamma)$ is a single essential simple closed curve, the
restriction $f:\delta\to\gamma$ has degree $2$, and the disk
$f^{-1}(D)$ bounded by $\delta$ contains the branch point $12/5$ and
exactly one point of $A$.
\end{lemma}
\begin{proof}
By Lemma~\ref{lem:HK-map}, the branch values of $f$ are exactly
$\infty$ and $v$.  The disk $D$ contains exactly one branch value,
namely $v$, and exactly one non-branch value $a$.  Therefore the total
ramification of $f$ over $D$ is $1$: the unique simple branch point
$12/5$ maps to $v$.

The monodromy of the degree-two cover $f$ around the boundary
$\partial D=\gamma$ is the product of the local monodromies about the
marked points lying in $D$.  Since $v$ is a branch value of $f$, the
local monodromy about $v$ is the nontrivial transposition of the two
sheets, whereas the local monodromy about $a$ is trivial.  Hence the
total monodromy is nontrivial, and $f^{-1}(\gamma)$ is connected: it is
a single simple closed curve $\delta$, and $f:\delta\to\gamma$ is a
degree-$2$ covering.

Applying the Riemann--Hurwitz formula to the proper degree-two map
$f:f^{-1}(D)\to D$ over the disk $D$, with one branch point of
multiplicity $2$ over $v$, gives
\[
\chi\!\left(f^{-1}(D)\right)=2\,\chi(D)-1=2\cdot1-1=1.
\]
Since $f^{-1}(D)$ is connected and has a single boundary component
$\delta$, it is a disk.

Finally, we determine the $A$-points lying in $f^{-1}(D)$.  Since
$D\cap B=\{a,v\}$ and $A\subset B$, any point of $A$ lying in
$f^{-1}(D)$ must map into $D\cap A=\{a\}$, so
\[
f^{-1}(D)\cap A\subset f^{-1}(a).
\]
By Lemma~\ref{lem:HK-map}, each fiber over $a\in\{0,1,\tfrac34\}$ has
exactly two points, one in $A$ and one not in $A$:
\[
f^{-1}(0)=\{\tfrac34,-2\},\qquad
f^{-1}(1)=\{\tfrac65,\infty\},\qquad
f^{-1}(\tfrac34)=\{1,-6\}.
\]
The point $12/5\in f^{-1}(D)$ is a branch point and lies over
$v\notin A$, so it is not an $A$-point.  Therefore
$f^{-1}(D)\cap A$ consists of the unique $A$-point inside the fiber
$f^{-1}(a)$; explicitly it is $\tfrac34$ for $a=0$, $\infty$ for
$a=1$, and $1$ for $a=\tfrac34$.
\end{proof}

\subsection{The squared twist is special liftable}

% RECORDS: 564
\begin{lemma}
\label{lem:lift-square}
Let $a\in\{0,1,\tfrac34\}$ and $\gamma=\gamma_{a,v}$.  Then
$T_{\gamma}^{2}$ is $f$-liftable, with a lift equal to the full twist
$T_{\delta}$ about $\delta=f^{-1}(\gamma)$.  Moreover
$\HF(T_{\gamma}^{2})=1\in\Mod_{A}$ and
$\FAB(T_{\gamma}^{2})=1\in\Mod_{A}$.
\end{lemma}
\begin{proof}
By Lemma~\ref{lem:disk-preimage}, $f^{-1}(\gamma)=\delta$ is a single
curve mapping to $\gamma$ with degree $2$.  Choose annular coordinates
$(r,\theta)$ near $\gamma$ and $(r,\theta')$ near $\delta$ in which the
covering is $(r,\theta')\mapsto(r,2\theta')$.  A representative of
$T_{\gamma}^{2}$ has the form
\[
(r,\theta)\longmapsto(r,\theta+4\pi\,\rho(r))
\]
for a cut-off function $\rho$ supported in the annulus and changing by
one full twist across it.  Its lift is
\[
(r,\theta')\longmapsto(r,\theta'+2\pi\,\rho(r)),
\]
which is precisely the full Dehn twist $T_{\delta}$ along $\delta$.
Hence $T_{\gamma}^{2}\in\Lf$, with $\HF$-image realised by $T_{\delta}$.

We now determine $\HF(T_{\gamma}^{2})=1$.  The lift
$T_{\delta}\in\Mod_{f^{-1}(B)}$ acts by twisting about $\delta$.  By
Lemma~\ref{lem:disk-preimage}, after forgetting the points of
$f^{-1}(B)\setminus A$, the curve $\delta$ bounds a disk containing
exactly one point of $A$.  Any homeomorphism of a disk fixing a single
interior marked point and equal to the identity on the boundary is
isotopic relative to the boundary and the marked point to the identity
(the pure mapping class group of a disk with one marked point is
trivial; this is the Alexander trick).  Hence the class of $T_{\delta}$
in $\Mod_{A}$ is trivial, and
\[
\HF(T_{\gamma}^{2})=1\in\Mod_{A}.
\]

For the forgetful image $\FAB(T_{\gamma}^{2})$, note that after
forgetting the point $v$ from the marked set, the curve $\gamma$
bounds a disk whose intersection with $A$ is exactly $\{a\}$.  By the
same Alexander-trick argument, $T_{\gamma}$ and hence $T_{\gamma}^{2}$
is trivial in $\Mod_{A}$:
\[
\FAB(T_{\gamma}^{2})=1\in\Mod_{A}.
\]
\end{proof}

% RECORDS: 564
\begin{theorem}
\label{thm:preservation}
For each $a\in\{0,1,\tfrac34\}$, let $\gamma_{a,v}$ be a simple closed
curve separating $\{a,v\}$ from the other three points of $B$.  Then
\[
T_{\gamma_{a,v}}^{2}\in\Sf,
\]
and consequently
\[
T_{\gamma_{a,v}}^{2}\cdot\Df=\Df.
\]
In particular, $\gamma=\gamma_{0,v}$ satisfies the preservation
requirement of the original problem with power $2$.
\end{theorem}
\begin{proof}
Let $\gamma=\gamma_{a,v}$.  By Lemma~\ref{lem:lift-square},
$T_{\gamma}^{2}\in\Lf$ and both $\HF(T_{\gamma}^{2})$ and
$\FAB(T_{\gamma}^{2})$ equal the identity in $\Mod_{A}$.  Hence
$T_{\gamma}^{2}\in\Sf=\{g\in\Lf:\HF(g)=\FAB(g)\}$.  Applying
Proposition~\ref{prop:HK21} with $g=T_{\gamma}^{2}\in\Sf$ gives
\[
T_{\gamma}^{2}\cdot\Df=\Df.
\]
Taking $a=0$ gives the explicit curve $\gamma_{0,v}$.
\end{proof}

\section{The vanishing-cycle twist lies in $\Ef$}
\label{sec:Ef}

We now turn to the action on components.  We will show that the
squared twist $T_{\gamma_{0,v}}^{2}$ lies in the subgroup $\Ef$ coming
from loops in the Hironaka--Koch parameter space $V$, and hence fixes
the base component of $\Df$.

% RECORDS: 643
\begin{lemma}
\label{lem:collision-in-V}
In the Hironaka--Koch parameter space $V=\Per_{4}(0)^{*}$, the collision
$v=c_{1}$ occurs as a removed puncture.  In the coordinates of
\cite[\S4.1]{HK}, with the periodic critical point normalized to
$c_{1}=0$, the collision condition is equivalent to the conic
\[
G:\;-x^{2}+y+2x-1=0.
\]
In those coordinates the one-parameter family $V$ is represented by the
diagonal $x=y$ with the collision and degeneracy loci removed, so
\[
G\cap V:\quad x^{2}-3x+1=0,
\]
and the collision punctures are
\[
x=\frac{3\pm\sqrt5}{2}.
\]
The associated vanishing cycle transported to the base fibre is
$\gamma_{0,v}$ under the Hironaka--Koch marking convention.
\end{lemma}
\begin{proof}
This is the coordinate computation of \cite[\S4.1]{HK}.  In those
coordinates the periodic critical point is normalized to $c_{1}=0$, and
the free critical value is a rational function $z(x,y)$.  Hironaka--Koch
write the collision condition as
\[
z(x,y)=0\quad\Longleftrightarrow\quad -x^{2}+y+2x-1=0,
\]
i.e.\ the collision conic is $y=(x-1)^{2}$.  The one-parameter family
$V$ is given in this chart by the diagonal $x=y$, after removing the
finite set on which the five marked points fail to be distinct or the
rational map degenerates.  Substituting $y=x$ into $G$ gives
\[
-x^{2}+x+2x-1=0\quad\Longleftrightarrow\quad x^{2}-3x+1=0,
\]
with roots $x=(3\pm\sqrt5)/2$.  At these parameters the two marked
points $0$ and $v$ coincide, so the marked sphere is not in
$\Teich(B)$; such parameters are punctures of $V$.  A small loop
around either puncture lies in $V$, and the colliding pair is
$\{0,v\}$; therefore the vanishing cycle transported to the base fibre
is, by definition of $\gamma_{0,v}$, the curve separating $\{0,v\}$
from the remaining three points of $B$.
\end{proof}

% RECORDS: 262,643
\begin{theorem}
\label{thm:Ef}
With the Hironaka--Koch base point $*_{V}$ and marking conventions, the
squared vanishing-cycle twist satisfies
\[
T_{\gamma_{0,v}}^{2}\in\Ef.
\]
\end{theorem}
\begin{proof}
Let $p$ denote one of the collision punctures
\[
x=\frac{3\pm\sqrt5}{2}
\]
from Lemma~\ref{lem:collision-in-V}, and let $U$ be a small coordinate
disk about $p$ inside the diagonal with local holomorphic coordinate
$t$ centred at $p$, so that $U\setminus\{p\}\subset V$.  The local
meridian $t=\varepsilon e^{2\pi i\theta}$, $0\le\theta\le1$, together
with the Hironaka--Koch transport path in $V$ from $U$ to the base
point $*_{V}$, defines a based loop $\ell$ in $V$.

The collision has multiplicity two.  This is the coordinate-invariant
content of Theorem~\ref{thm:normal-form}: in a one-parameter model of
the collision, the relative coordinate between the periodic critical
point and the free critical value is a nonzero holomorphic function
times $t^{2}$.  The normal-form family of Theorem~\ref{thm:normal-form}
and the diagonal in the Hironaka--Koch chart are two holomorphic
parametrizations of the same one-parameter locus $\Per_{4}(0)$ near the
collision.  The vanishing order of the relative coordinate of two marked
points at a puncture is invariant under holomorphic reparametrization of
the parameter and under M\"obius changes of coordinate on the sphere.
Hence the order-two vanishing proved in Theorem~\ref{thm:normal-form}
also holds for the local parameter $t$ on the Hironaka--Koch diagonal.
After possibly shrinking $U$, we may therefore write
\[
w(t)=u(t)\,t^{2},\qquad u(0)\neq0,
\]
for the relative coordinate of the colliding marked points $0$ and $v$.
As $t$ goes once around $0$, $w(t)$ winds twice around $0$.

Inside a small disk containing only the two colliding marked points, the
resulting motion is a pure braid with relative winding number two; by
the standard braid-to-mapping-class identification in a twice-punctured
disk (Farb--Margalit \cite[Chapter~9, Theorem~9.1 and \S9.1]{FM}), this
is the square of the Dehn twist about the curve surrounding the two
punctures, up to orientation.  The local vanishing cycle transported to
the base fibre is $\gamma_{0,v}$ by
Lemma~\ref{lem:collision-in-V}.

The exact identification of the based meridian's monodromy uses the
Hironaka--Koch base point and marking conventions: the base fibre over
$*_{V}$ is identified with $(\Sph^{2},B)$, the colliding pair is
$\{0,v\}$, and the transport path from $*_{V}$ to the puncture is the
one used in that marking.  Under this convention the vanishing cycle
transported to the base fibre is the specific curve $\gamma_{0,v}$.
Thus the based monodromy of $\ell$ is $T_{\gamma_{0,v}}^{\pm2}$; since
$\Ef$ is a subgroup, this implies $T_{\gamma_{0,v}}^{2}\in\Ef$.
\end{proof}

% RECORDS: 643
\begin{corollary}
\label{cor:base-fixed}
Let $D_{0}$ be the connected component of $\Df$ containing the base
point $\tau_{*}$.  Then
\[
T_{\gamma_{0,v}}^{2}(D_{0})=D_{0}.
\]
\end{corollary}
\begin{proof}
Combine Theorem~\ref{thm:Ef} with Lemma~\ref{lem:base-component}.
\end{proof}

\begin{remark}
\label{rem:caveat}
The arguments above do not show that $\Ef$ is normal in $\Sf$.
Therefore Corollary~\ref{cor:base-fixed} rigorously fixes only the
base component.  It does not by itself determine the action of
$T_{\gamma_{0,v}}^{2}$ on all components of $\Df$, and it does not rule
out nontrivial behaviour on other components.  However, no argument in
this paper establishes an infinite orbit for this Dehn twist.
\end{remark}

\section{Consequences and open points}
\label{sec:open}

The original problem asks for an explicit simple closed curve $\gamma$
such that a power of $T_{\gamma}$ preserves $\Df$ and acts on
$\pi_{0}(\Df)$ with an infinite orbit.  The results above give the
following partial answer.
The preservation requirement is solved by the explicit curve
\[
\gamma=\gamma_{0,v}
\]
with power $2$: Theorem~\ref{thm:preservation} proves
$T_{\gamma_{0,v}}^{2}\in\Sf$, and Proposition~\ref{prop:HK21} then gives
\[
T_{\gamma_{0,v}}^{2}\cdot\Df=\Df.
\]
More generally, for each $a\in\{0,1,\tfrac34\}$, the squared twist about
$\gamma_{a,v}$ lies in $\Sf$ and preserves $\Df$.

The infinite-orbit requirement is not proved for this Dehn twist.
Theorem~\ref{thm:Ef} shows
\[
T_{\gamma_{0,v}}^{2}\in\Ef,
\]
so the base component $D_{0}$ of $\Df$ is fixed by
Corollary~\ref{cor:base-fixed}.  The known Hironaka--Koch proof that
$\Df$ has infinitely many components uses the braid element
\[
g=s_{2}\,c\,r_{1}\,s_{1}\,c\,r_{2},
\]
whose image under the Hironaka--Koch quotient map has infinite order;
see \cite[Theorem~1.3 and Remark~4.7]{HK}.  The arguments of this paper
do not identify this element, or any power of it, with a Dehn twist
about a single simple closed curve on $(\Sph^{2},B)$.  The infinite
orbit produced by $g$ is therefore not realised here by a Dehn twist.

The precise open points are the following.
\begin{enumerate}[label=(\arabic*)]
\item No argument in this paper shows that a power of $T_{\gamma}$ for
some single essential simple closed curve
$\gamma\subset(\Sph^{2},B)$ lies in $\Sf\setminus\Ef$, which is the
standard route in the Hironaka--Koch framework to obtaining nontrivial
component orbits.
\item The canonical curve $\gamma_{0,v}$ does not supply the infinite
orbit: its square lies in $\Ef$ and fixes the base component $D_{0}$.
\item Since $\Ef$ is not known here to be normal in $\Sf$, membership of
$T_{\gamma_{0,v}}^{2}$ in $\Ef$ controls the base component only.
Ruling out nontrivial orbits on all components would require additional
information, such as normality of $\Ef$ or a full coset calculation.
\item For the other squared $v$-pair twists
$T_{\gamma_{1,v}}^{2}$ and $T_{\gamma_{3/4,v}}^{2}$, preservation is
proved by Theorem~\ref{thm:preservation}, but their membership or
non-membership in $\Ef$ is not settled here.
\item The Hironaka--Koch infinite-component witness
$g=s_{2}cr_{1}s_{1}cr_{2}$ is a braid element; identifying it, or a
power of it, with a Dehn twist about an explicit simple closed curve
remains open.
\item A complete classification of essential simple closed curves whose
Dehn-twist powers lie in $\Sf\setminus\Ef$ is not given here.
\end{enumerate}

\begin{thebibliography}{9}

\bibitem{HK}
E.~Hironaka and S.~Koch,
\emph{A disconnected deformation space of rational maps},
arXiv:1602.07378, 2016.

\bibitem{FM}
B.~Farb and D.~Margalit,
\emph{A Primer on Mapping Class Groups},
Princeton Mathematical Series, Princeton University Press, 2012.

\end{thebibliography}

\end{document}